\begin{proof}
\textbf{Trivial (degenerate) cases.}
\begin{enumerate}
\item \textbf{Case \(n=0\).}  
\[
(1+x)^{0}=1 \qquad\text{and}\qquad 1+0\cdot x =1 .
\]  
Equality holds.

\item \textbf{Case \(n=1\).}  
\[
(1+x)^{1}=1+x = 1+1\cdot x .
\]  
Equality holds.

\item \textbf{Case \(x=0\).} For any \(n\ge 0\),
\[
(1+0)^{n}=1 = 1+n\cdot 0 .
\]  
Equality holds.
\end{enumerate}
These three situations constitute all edge cases that will be excluded from the induction argument that follows.

\medskip
\textbf{Base of induction (\(n=2\)).}  
Using the binomial theorem (or direct multiplication) we have the algebraic identity
\begin{align*}
(1+x)^{2}=1+2x+x^{2},
\end{align*}
and since \(x^{2}\ge 0\) for every real \(x\),
\begin{align*}
(1+x)^{2}=1+2x+x^{2}\;\ge\;1+2x .
\end{align*}
Thus the inequality holds for \(n=2\).

\medskip
\textbf{Inductive step.}

\textit{Inductive hypothesis.}  
Assume that for some integer \(k\ge 2\) and for all \(x\ge -1\),
\begin{equation}
(1+x)^{k}\;\ge\;1+kx .\tag{IH}
\end{equation}

\textit{Multiplication by a non‑negative factor.}  
Because \(x\ge -1\) we have \(1+x\ge 0\). Multiplying both sides of \eqref{IH} by \(1+x\) preserves the inequality:
\begin{align}
(1+x)^{k+1}&=(1+x)^{k}(1+x) \nonumber\\
&\ge (1+kx)(1+x). \tag{1}
\end{align}

\textit{Expansion of the right‑hand side.}  
The product expands as
\begin{align}
(1+kx)(1+x)=1+(k+1)x+kx^{2}\;\ge\;1+(k+1)x .\tag{2}
\end{align}
Since \(k\ge 0\) and \(x^{2}\ge 0\), the term \(kx^{2}\) is non‑negative, giving the last inequality.

Combining \eqref{1} and \eqref{2} yields
\[
(1+x)^{k+1}\;\ge\;1+(k+1)x .
\]
Thus the statement holds for \(k+1\). By mathematical induction the inequality is true for every integer \(n\ge 2\).

Together with the trivial cases handled in Section 1, the inequality holds for **all** integers \(n\ge 0\) and all real numbers \(x\ge -1\).

\medskip
\textbf{Equality conditions.}  
Equality in the inductive step requires simultaneously  

\begin{enumerate}
\item equality in the hypothesis \((1+x)^{k}=1+kx\); and
\item vanishing of the non‑negative term \(kx^{2}\).
\end{enumerate}
Because \(k\ge 0\), the second condition forces either \(k=0\) (i.e. \(n=0\)) or \(x=0\).
The case \(n=1\) gives equality for every admissible \(x\) because the inequality reduces to the identity \(1+x=1+x\).
The boundary point \(x=-1\) yields \((1-1)^{n}=0\) and \(1-n\); equality occurs only when \(n=1\), already included.

Hence **equality holds precisely when**
\[
n=0\quad\text{or}\quad n=1\quad\text{or}\quad x=0 .
\]

\medskip
\textbf{Conclusion.}  
For every real number \(x\ge -1\) and every integer \(n\ge 0\),
\[
\boxed{\, (1+x)^{n}\;\ge\;1+nx \,}
\]
with equality exactly in the cases \(n=0\), \(n=1\) or \(x=0\). \qed
\end{proof}